\subsubsection{“复”对称性破缺}\label{compsym}



如果我们让 $\vec\phi$ 的所有分量在无穷远处都取期望值，同时要求 $A$ 在无穷远处为零，那么复联络 $\A$ 在 $y\to\infty$ 时为常值且对角，
\begin{align}\label{const}
\A_z &\sim \frac{1}{\zeta} \begin{pmatrix} \ca & 0 \cr 0 & -\ca \end{pmatrix} \cr
\A_{\bar z} &\sim \zeta \begin{pmatrix} \bar \ca & 0 \cr 0 & - \bar \ca \end{pmatrix} \cr
\A_y &\sim \begin{pmatrix} \ca_1 & 0 \cr 0 & -\ca_1 \end{pmatrix} ,\cr
\end{align} 其中 $\ca$ 为复数。$\zeta$ 的各次幂来自从 $\varphi$ 到 $\A$ 的变量替换。
% [Tutor Note: 约定差异] 这里 $\A_z$ 带因子 $1/\zeta$、$\A_{\bar z}$ 带因子 $\zeta$，这来自 $\A$ 的分解约定（$\A=\A_z\d z+\A_{\bar z}\d\bar z$ 与复化后的 $\varphi$ 的关系）。若采用相反的约定，(\ref{tomogro2}) 中 $v$ 的常数项 $\ca/\zeta$ 的符号会随之改变，核对时请以原文的这一组式子为准。

在 $z$ 方向的无穷远处，对任意 $y$，解都化为 Nahm 方程的唯一解（文献~\cite{Kronheimer}），该解在 $y=0$ 处有 Nahm 极点，并在 $y=\infty$ 时呈现 (\ref{const}) 的行为。特别地，当 $z\to\infty$ 而 $y$ 固定时，联络形式为常数。对于这种形式的联络，我们可以把小的全纯截面 $s$ 写成
\begin{equation}\label{goodform}
s = \begin{pmatrix} e^{- \zeta \bar \ca'\, \bar z} & 0 \cr 0 &
e^{ \zeta \bar \ca'\, \bar z}  \end{pmatrix} s_0
\end{equation}
其中 $\ca'$（在 $y\to\infty$ 时等于 $\ca$）是 $\A_{\bar z}/\zeta$ 在 $z$ 很大而 $y$ 固定时的常值，而 $s_0$ 在 $z$ 很大、$y$ 固定时有有限的极限。

为分析这类解中的全纯数据，把联络规范掉是不自然的。我们所需要的信息会隐藏在无穷远处所需规范变换的行为之中。因此，我们将选取一个复规范变换，把联络精确地（而不只是渐近地）变为
\begin{align}\label{consto}
\A_z &=\frac{1}{\zeta} \begin{pmatrix} \ca & 0 \cr 0 & -\ca \end{pmatrix} \cr
\A_{\bar z} &= \zeta \begin{pmatrix} \bar \ca & 0 \cr 0 & - \bar \ca \end{pmatrix} \cr
\A_y &= \begin{pmatrix} \ca_1 & 0 \cr 0 & -\ca_1 \end{pmatrix} \cr
\end{align}
正如我们马上会看到的，这可以由相当简单的一类规范变换做到。考虑 $\R^2\times \R_+$ 上的任一规范场，例如具有上述渐近形式、且当 $z\to\infty$（$y$ 固定）时联络形式趋于一个非零常数的规范场。若把 $\R^2$ 紧化为 $\CP^1$，我们会说这样的联络在 $z=\infty$ 处有一个带二阶极点的非正则奇点。具有非正则奇点的联络会出现 Stokes 现象。若只考虑没有本质奇点（即在无穷远处只有多项式增长）的规范变换，则 Stokes 数据是规范不变的。两个具有非正则奇点的平坦联络，当且仅当具有相同的单值性和相同的 Stokes 数据时，才可由一个在无穷远处只有多项式增长的规范变换彼此规范等价。

由于我们的联络在 $z=\infty$ 处只有一个二阶极点，它只有两个 Stokes 扇区。Stokes 理论告诉我们，绕 $z=\infty$ 的单值性可以分解为一个对角的形式单值矩阵与一串 Stokes 矩阵的乘积，后者交替为上三角与下三角矩阵，且对角元为 1。只有两个 Stokes 扇区时，单值性的表达式为
\begin{equation}\label{monodr}M=
\begin{pmatrix}1 & b \cr 0 & 1 \end{pmatrix}\begin{pmatrix}1 & 0 \cr \tilde b & 1 \end{pmatrix}\begin{pmatrix}\mu & 0 \cr 0 & \mu^{-1} \end{pmatrix} ,\end{equation}
其中 $b,\tilde b,$ 与 $\mu$ 为常数。由于我们在单连通的 $\R^2$ 上，无穷远处的单值性必须是 $M=1$。连同单值性的形式 (\ref{monodr})，这意味着 (\ref{monodr}) 中的三个因子——两个 Stokes 矩阵与形式单值性——都必须等于 1。因此，我们可以用一个在无穷远处只有多项式增长的规范变换，把联络变到常值对角形式 (\ref{consto})。
% [Tutor Note: 省略证明] “三个因子都必须等于 1”并非显然：把 (\ref{monodr}) 乘开可得 $M_{11}=(1+b\tilde b)\mu$、$M_{12}=b\mu^{-1}$、$M_{21}=\tilde b\mu$、$M_{22}=\mu^{-1}$。于是 $M=1$ 先给出 $\mu=1$，再由 $M_{12}=b=0$、$M_{21}=\tilde b=0$ 定出两个 Stokes 参数为零；读者不妨自己乘一遍以确认。

特别地，$s$ 现在将取形式
\begin{equation}
s = \begin{pmatrix} e^{- \zeta \bar \ca\, \bar z} & 0 \cr 0 & e^{ \zeta \bar \ca\, \bar z}  \end{pmatrix} \begin{pmatrix}P(z) \cr Q(z) \end{pmatrix}
\end{equation}
其中 $P$ 与 $Q$ 为多项式。方程 $s \wedge \D_z s=K(z)$ 给出
\begin{equation}\label{normo}
\left(P \partial_z Q - Q \partial_z P\right) - \frac{2 \ca P Q}{\zeta} = K(z)
\end{equation}
它可以化为要求 $Ke^{2 \ca z}/Q^2 $ 在 $Q(z)$ 的零点 $w_i$ 处没有留数。这变成
\begin{equation}\label{bethetwist}
\frac{2 \ca}{\zeta} + \sum_a \frac{k_a}{w_i - z_a} = \sum_j \frac{2}{w_i - w_j}
\end{equation}
% [Tutor Note: 易错点] (\ref{bethetwist}) 右端写作对全部 $j$ 求和，$j=i$ 项本身是极点，应理解为去掉该极点后的正则部分（即 $j\ne i$ 之和）；与 (\ref{bethe}) 比较时注意：取 $\ca=0$ 后二者一致，而 $\ca\ne0$ 时指数因子 $e^{2\ca z}$ 把二阶极点处的留数条件平移了 $\frac{2\ca}{\zeta}$。
我们所关心的是 (\ref{normo}) 的解模去标度变换 $P\to\lambda P$、$Q\to
\lambda^{-1}Q$（它对应对角平坦联络的一个自同构），因此实际上我们只关心 $Q(z)$ 的零点。

这些仍然是 Bethe 方程，只不过这次是 $z=\infty$ 处有非正则奇点的 Gaudin 模型的 Bethe 方程（文献~\cite{Ry,FFT,FFRy}）。文献~\cite{FFRy} 证明，这些 Bethe 方程的解与自旋链的本征向量一一对应，总共有 $\prod_a(k_a+1)$ 个解。也许也可以从 XXX 自旋链的理论中得出这一结果，因为 XXX 链比 Gaudin 模型更为人们所熟悉。Gaudin 模型的 Bethe 方程来自 XXX 链 Bethe 方程的某个特定的“大杂质极限”。形变参数 $\ca$ 对应 XXX 链的扭曲；众所周知，该扭曲把整体 $SU(2)$ 对称性破缺为 $U(1)$，并简化 Bethe 方程解的计数：不再是每个作为 $SU(2)$ 最高权的自旋链哈密顿量本征向量各对应一个解，而是每个本征向量都对应一个解。

作为一个简单例子，考虑单个荷为 2 的 't Hooft 算符，即 $K(z)=z^2$。正如第~\ref{physhilb}~节那样，(\ref{normo}) 有一个 $P=1$、$Q$ 为 $z$ 的二次多项式的解，还有一个 $Q=1$、$P$ 为二次多项式的解。若 $P$ 与 $Q$ 都是 $z$ 的一次式，会怎样呢？同样如第~\ref{physhilb}~节，我们可以用 $P=Q=z\sqrt{-\zeta/2\ca}$ 求解 (\ref{normo})，但这个解不对应 oper，因为 $P$ 与 $Q$ 有公共零点。新鲜之处在于，还存在一个可接受的解 $P=-(\zeta/2\ca)(z-\zeta/\ca)$、$Q= z+\zeta/\ca$，它对应于 Bethe 方程 (\ref{bethetwist}) 的一个单 Bethe 根解。这样总共有 $2+1=3$ 个满足必要条件的 oper。最后一个解在取 $\ca=0$ 时消失，而若取 $\zeta\to 0$，它则使 $P$ 与 $Q$ 出现公共零点。



\subsection{Oper 与应力张量}\label{opstress}

在我们对平凡单值性 oper 的分析中，我们使用了联络平凡的规范，即 $\A_z=\A_{\bar z}=0$。相应地，我们不得不为生成全纯子丛 $\L\subset E$ 的小截面 $s$ 作一般性的拟设：
\begin{equation}\label{poxo} s=\begin{pmatrix}P\cr Q\end{pmatrix}.\end{equation}
这里我们将作一个规范变换，把 $s$ 变为上分量为零的标准形式，并看看关于 $\A_z$ 能说些什么。这样做有两点好处。我们将开始理解 oper 与共形场论的关系；而且我们将得到一种更为一般的描述，不再局限于平凡单值性的 oper 丛这一情形（不过该情形在本文中尤为重要）。

用光滑规范变换把 $s$ 变为简单形式会使 $\A_z$ 与 $\A_{\bar z}$ 都不为零。事实证明，保持 $\A_{\bar z}=0$ 更有帮助。若还要求 $\A_z$ 正则，就需要我们的规范变换是全纯的，而这个条件过于苛刻。因此我们转而考虑亚纯规范变换，它可以保持 $\A_{\bar z}=0$、把 $s$ 变为标准形式，并在 $\A_z$ 中产生极点。

把 $s$ 变为标准形式的最显然的亚纯规范变换是
\begin{equation}\label{helm} h=\begin{pmatrix}Q & -P\cr 0 & Q^{-1}\end{pmatrix}.\end{equation}
我们有
\begin{align}\label{elm}  hs& = \begin{pmatrix}0 \cr 1\end{pmatrix} \cr
          h\partial_z h^{-1} & = \partial_z+
          \begin{pmatrix} -\frac{Q'}{Q} & -K  \cr 0 & \frac{Q'}{Q}\end{pmatrix},\end{align}
其中一如前述 $K=PQ'-QP'$。事实证明，更方便的是转到 $\A_z$ 的右上矩阵元为 $-1$ 的规范，代价是把 $s$ 变为自身的倍数（这不改变由 $s$ 生成的线丛）。我们再作一次规范变换
\begin{equation}\label{furth}\tilde h=
\begin{pmatrix}1/\sqrt{K} & 0 \cr 0 & \sqrt{K}\end{pmatrix}.\end{equation}
$\sqrt{K}$ 可能出现的双值性无关紧要，理由如下。若 $G^\vee=SO(3)$，从而 $G=SU(2)$，则 $K$ 为完全平方，因其所有零点都是偶次的。反之，若 $G^\vee=SU(2)$，则 $G=SO(3)$，而我们其实应当用伴随表示而非二维表示来写所有公式；相应地，规范变换的符号无关紧要。经过 $\tilde h$ 的规范变换后，$\A_z$ 取形式
\begin{equation}\label{omogro}\A_z=
\begin{pmatrix} -v & -1\cr 0 & v \cr\end{pmatrix}\end{equation}
其中我们已令
\begin{equation}\label{tomogro}v=-\frac{ K'}{2K}+\frac{Q'}{Q}=-\sum_a \frac{ k_a/2}
{z-z_a}+\sum_i\frac{1}{z-w_i}.\end{equation}
在最后一步中我们用了 $K=\prod_a (z-z_a)^{k_a}$、$Q=\prod_i(z-w_i)$。最后，一个下三角规范变换
\begin{equation}\label{morgo} \begin{pmatrix} 1 & 0 \cr v & 1\end{pmatrix}\end{equation}
给出 $\A_z$ 的最终结果：
\begin{equation}\label{torog} \A_z=\begin{pmatrix} 0 & -1 \cr \bt & 0 \end{pmatrix},\end{equation}
其中
\begin{equation}\label{orog} t= -v'-v^2.\end{equation}
一般地，$\bt$ 在 $z_a$ 与 $w_i$ 处都有极点。在 $z=z_a$ 附近，
\begin{equation}\label{roog}t\sim-\frac{ j_a(j_a+1)}{(z-z_a)^2}+\frac{c_a}{z-z_a}+\dots, ~ j_a=k_a/2.\end{equation}
在 $z=w_i$ 附近，
\begin{equation}\label{frog}t\sim  \frac{1}{z-w_i}
\left(\sum_a\frac{k_a}{w_i-z_a}-\sum_{j\not=i}\frac{2}{w_i-w_j}\right).\end{equation}
于是，$\bt$ 在 $z=w_i$ 处没有奇点，当且仅当 Bethe 方程 (\ref{bethe}) 得到满足。

为获得更多理解，令 $\D_z=\partial_z+[\A_z,\,\cdot\,]$，并寻找一个平坦截面，即 $\D_z\begin{pmatrix}f\cr \tilde f\end{pmatrix}=0$ 的全纯解。我们发现 $\tilde f=f'$，且
\begin{equation}\label{opeq}\left(\frac{\partial^2}{\partial z^2} +t(z)\right)f = 0.\end{equation}
我们是在某个特定的局部坐标 $z$ 下完成这一推导的，但我们出发时的概念——带 oper 结构的平坦丛——并不依赖于局部坐标。因此方程 (\ref{opeq}) 在局部坐标变换下必须是协变的，$\bt$ 也须有相应的变换。一个简短的计算（或对上述推导更仔细的研究）表明，在局部坐标变换 $z\to \tilde z$ 以及 $\bt$ 的相应变换下，$\bt$ 这个对象的变换方式如同二维共形场论中的应力张量。换言之，它并不像人们可能从它与 (\ref{opeq}) 中二阶导数 $\partial^2/\partial z^2$ 的配对出发所朴素设想的那样按二次微分变换，而是带有一个涉及 Schwarzian 导数 $\frac{1}{2} \{z,\tilde z\}$ 的“异常”项。$t$ 在 $z=z_a$ 处的二阶极点就好像在 $z_a$ 处插入了一个初级场。
% [Tutor Note: 易错点] $\bt$ 变换中的 Schwarzian 系数是 $\frac12$，不是 $\frac1{12}$：后者属于量子应力张量 $T$ 的反常项 $\frac{c}{12}\{z,\tilde z\}$；用 $\bt=b^2T$ 与 $c=1+6(b+b^{-1})^2\approx 6/b^2$ 换算，$b^2\cdot\frac{c}{12}\to\frac12$，两者一致。自己推 Schwarzian 项时最容易在这里丢因子。


如果我们从 $\A_z = \zeta^{-1} \mathrm{diag}(\ca,-\ca)$ 出发，并作同一串规范变换，其中 $K=PQ'-QP' - 2 \ca \zeta^{-1} PQ$，那么我们会得到相同的公式，只是 $v$ 中多出一个常数项：
\begin{equation}\label{tomogro2}v=- \frac{\ca}{\zeta}-\sum_a \frac{ k_a/2}
{z-z_a}+\sum_i\frac{1}{z-w_i}.\end{equation}
现在 $\bt(z)$ 在 $z=\infty$ 处有一个四阶极点。这在共形场论中将对应于在无穷远处插入一个略不寻常的算符（文献~\cite{G}）。


最后，Bethe 方程具有下列有趣性质：在 $z_a$（以及 $\ca$）固定时，它们描述一个 Yang-Yang 函数的驻点：
\begin{align}\label{yy}
\W(w_i,z_a) = -  \sum_{i < j} \log((w_i - w_j)^2) + \sum_{i,a} k_a \log(w_i - z_a) & -
 \frac{1}{4} \sum_{a < b} k_a k_b \log( (z_a - z_b)^2) \cr & +\frac{2 \ca}{\zeta} \sum_i w_i - \frac{\ca}{\zeta} \sum_a k_az_a
\end{align}
换言之，Bethe 方程可以写成
\begin{equation}
\frac{\partial \W}{\partial w_i} =0.
\end{equation}
% [Tutor Note: 验证计算] 对 (\ref{yy}) 求 $\partial_{w_i}$：第一项中每一项 $\log((w_i-w_j)^2)$ 对 $w_i$ 的导数为 $2/(w_i-w_j)$，故第一项共给出 $-\sum_{j\ne i}2/(w_i-w_j)$；第二项给出 $\sum_a k_a/(w_i-z_a)$；与 $\ca$ 成正比的线性项给出 $2\ca/\zeta$。三者相加为零正是 (\ref{bethetwist})。这也解释了“Yang-Yang 函数”之名：Bethe 根恰好是它的驻点。

Yang-Yang 函数还有另一个有趣性质：$\bt$ 在 $z=z_a$ 处单极点的系数——有时称为 accessory 参数——由下式给出：
\begin{equation}
c_a = \frac{\partial \W}{\partial z_a}.
\end{equation}
$\W$ 中含有一些与 $w_i$ 无关的项，它们被引入正是为了保证这一关系成立。类似的关系对 $\ca \partial_\ca \W$ 也成立。

\subsection{Oper 与 Virasoro 共形块}\label{opbl}



本节我们将说明这些公式如何在 Virasoro 共形块的半经典极限中自然地出现。半经典极限定义为这样的极限：Virasoro 代数的中心荷 $c$ 趋于无穷，同时各算符的共形维数也按与 $c$ 相同的方式标度。

中心荷有一种有用的参数化方式：$c=1+6 Q^2$，其中 $Q=b + b^{-1}$。在此参数化下，我们取 $b \to 0$ 以得到半经典极限。算符的共形维数可方便地参数化为 $\Delta = \alpha(Q-\alpha)$。参数 $\alpha$ 常被称为“动量”。当 $b\to 0$ 时，我们保持 $b \alpha=\eta$ 固定。于是，在关联函数中插入能动量张量 $T$ 的贡献按 $b^{-2}$ 标度，我们可以定义有限极限 $\bt = b^2 T$。我们建议把这个插入到某些共形块中的 $\bt$ 与第~\ref{opstress}~节的 $\bt$ 等同起来。

量子应力张量 $T(z)$ 在共形变换下有反常；它平移 Schwarzian 导数的一个倍数 $\frac{c}{12} \{z,\tilde z\}$。因此 $\bt=b^2T$ 具有一个在 $b\to 0$ 时不依赖于 $b$ 的共形反常。上一节中 $\bt$ 在 $z=z_a$ 附近的行为，对应于 $\alpha_a = - \frac{k_a}{2b}$ 的 Virasoro 初级场附近的行为。这类算符非常特殊：只涉及这类算符的关联函数与共形块，可以很容易地用自由场实现来描述。我们可以仿照上一节的讨论来描述这样的实现。

令 $\chi$ 为二维自由场，其两点函数为 $\langle\chi(z)\chi(z')\rangle=-\frac{1}{2} \ln(z-z')$。构造中心荷 $c=1+6 Q^2$ 的能动量张量的标准做法是取
\begin{equation}
T = - :\partial \chi \partial\chi: + Q \partial^2 \chi.
\end{equation}
若定义 $v=-b \partial \chi$，则该定义在 $b \to 0$ 的极限下化为 (\ref{orog})。这是自由场实现可能对我们有用的第一个迹象。维数为 $\Delta = \alpha(Q-\alpha)$ 的算符可以很容易地用自由玻色子的正规序指数来描述，
\begin{equation}
V_{\alpha}(z) = :e^{2\alpha \chi(z)}:\, .
\end{equation}

第二个迹象来自 (\ref{tomogro})：量 $v$ 是在 $z=z_a$ 处动量为 $- {k_a}/{2b}$、在 $z=w_i$ 处动量为 $1/b$ 的手征顶点算符存在时，$-b \partial \chi$ 期望值的半经典极限。算符 $V_{1/b}(w_i)$ 的维数为 $1$。在自由场实现中它们通常被称为 screening 算符（文献~\cite{DF,Felder}），并且自然地在曲线上积分。应力张量在 screening 算符位置附近的奇异部分
\begin{equation}
T(z) V_{1/b}(w_i) \sim \frac{1}{(z-w_i)^2} V_{1/b}(w_i)  +
\frac{1}{z-w_i} \partial_{w_i} V_{1/b}(w_i) +\cdots = \partial_{w_i}
\left( \frac{1}{z-w_i} V_{1/b}(w_i) \right) \cdots
\end{equation}
是全导数，因而在 screening 算符的位置被积分掉后便消失。
% [Tutor Note: 物理图像] 这一句是整套自由场实现的关键：screening 算符插入点的积分因上述全导数性质而不依赖于 $w_i$ 的位置，于是积分结果自动满足 Virasoro Ward 恒等式，即得到一个共形块；而半经典极限下该积分由 Yang-Yang 函数的鞍点（即 Bethe 根）主导，这正是前文 oper 数据与共形块之间的桥梁。

我们可以很容易地算出如下自由场关联函数：
\begin{align}\label{easy}
\left\langle \prod_i  V_{1/b}(w_i)
\prod_a  V_{-k_a/2b}(z_a)  \right\rangle_{\mathrm{free}}& = \prod_{i < j}(w_i - w_j)^{-\frac{2}{b^2}
}\prod_{i,a}(w_i - z_a)^{\frac{k_a}{b^2}}
\prod_{a <b}(z_a - z_b)^{-\frac{1}{2b^2}k_a k_b}\cr & = \exp\left({\frac{1}{b^2} \W(w,z)}\right).
\end{align}
右端的指数正是 Yang-Yang 函数 (\ref{yy})！（眼下与 $\ca$ 成正比的项不出现，因为我们尚未引入对称性破缺。）

考虑积分
\begin{equation} \label{freeinto}
\left\langle  \prod_a  V_{-k_a/2b}(z_a)
\right\rangle_{\Gamma} = \int_\Gamma \left\langle \prod_i
 V_{1/b}(w_i)  \prod_a  V_{-k_a/2b}(z_a)  \right\rangle_{\mathrm{free}}
\prod_i \d w_i,
\end{equation}
其中 $\Gamma$ 是使积分收敛的任一积分链。由 (\ref{easy})，这等价于
\begin{equation} \label{freeint}\left
\langle  \prod_a  V_{-k_a/2b}(z_a)\right
\rangle_{\Gamma} = \int_\Gamma \,\,\exp(\W(z_a,w_i)/b^2)\,\,\prod_i \d w_i.
\end{equation}


一个初级场乘积期望值的 Virasoro 共形块，是与 Virasoro Ward 恒等式相容的候选关联函数：
\begin{equation}\left
\langle  T(z) \prod_a  V_{-k_a/2b}(z_a)  \right\rangle_{\Gamma} =
\left( \sum_a \frac{\Delta_a}{(z-z_a)^2} +
\frac{1}{z-z_a} \frac{\partial}{\partial z_a} \right)\left \langle
\prod_a  V_{-k_a/2b}(z_a) \right \rangle_{\Gamma}
\end{equation}
由 (\ref{freeinto}) 或 (\ref{freeint}) 定义的函数，对 $w$ 的数目与积分链 $\Gamma$ 的任意选取都具有这一性质；这是利用定义 (\ref{freeinto}) 以及 screening 荷是维数为 1 的初级场这一事实来证明的。（函数 $\left\langle  T(z) \prod_a  V_{-k_a/2b}(z_a) \right \rangle_{\Gamma} $ 由积分 (\ref{freeinto}) 定义，并在右端的自由场关联函数中插入 $T(z)$。）我们刚才描述的正是共形块的自由场实现（文献~\cite{DF,Felder}）。

当点 $z_a$（$a=1,\dots,d$）变动时，积分 (\ref{freeint}) 所有可能的积分链 $\Gamma$ 构成的空间带有一个自然的平坦联络（Gauss-Manin 联络）。因此，函数 $\langle  \prod_a  V_{-k_a/2b}(z_a)
\rangle_\Gamma$——对任何固定的 $w$ 数目——给出 $d$ 股辫群的一个表示。已知（文献~\cite{Lawrence,SV}）这些正是与 Jones 多项式及其推广相联系的辫群表示。Jones 多项式与共形场论的这一关系，比第~\ref{relconf}~节中提到的经由 Knizhnik-Zamolodchikov 方程的关系，对本文更为有用。

函数 $\langle  \prod_a  V_{-k_a/2b}(z_a)
\rangle_\Gamma$ 并非维数与 $V_{-k_a/2b}$ 相同的初级场乘积的全部可能共形块；更确切地说，若算符 $V_{-k_a/2b}$ 是文献~\cite{BPZ} 所引入意义下的退化初级场，则它们就是全部这样的共形块。另一种说法是：若由 $\bt=b^2T$ 的小 $b$ 极限得到的 oper 在 $z=z_a$ 处被要求具有平凡单值性，则这些函数就是全部可能的共形块。这些概念以及它们之间的关系将在第~\ref{degtriv}~节中描述。

\subsubsection{非正则情形}\label{irreg}

现在让我们把复对称性破缺参数 $\ca$ 纳入上述讨论。我们当然可以在对 $w$ 的积分 (\ref{freeint}) 中修改指数，使其包含 Yang-Yang 函数 (\ref{yy}) 中与 $\ca$ 成正比的项。但这在共形场论中意味着什么呢？我们需要把自由场关联函数 (\ref{easy}) 换成
\begin{align}\label{oddball}
&\left\langle \prod_i
V_{1/b}(w_i)  \prod_a   V_{-k_a/2b}(z_a) \right \rangle_{\mathrm{free}}  \cr & =
\prod_{i \neq j}(w_i - w_j)^{-\frac{2}{b^2}}\prod_{i,a}(w_i - z_a)^{\frac{k_a}{b^2}}
\prod_{a \neq b}(z_a - z_b)^{-\frac{1}{b^2}k_a k_b}
e^{\frac{\ca}{\zeta b^2}
 \left( 2 \sum_i w_i - \sum_a z_a  \right)}.
 \end{align}
这一公式的共形场论解释是什么？

差不多从构造上就能看出，右端是给定场乘积的自由场关联函数，只是 $\chi$ 在无穷远处具有特殊的边界条件，$\chi \sim {\ca z}/{b \zeta} $。另一种说法是，我们在无穷远处插入了一个“非正则顶点算符”，即下式的 $L \to \infty$ 极限：
\begin{equation}
\exp\left({-\frac{2\ca}{\zeta b} L^2 \partial \chi(L) }\right).
\end{equation}


存在这样的非正则顶点算符时，应力张量具有预期的四阶极点。上述公式除了一点之外没有变化，即可能的积分围道的选择变多了。积分的结果是一个在无穷远处带有非正则穿刺点的共形块（其定义见文献~\cite{G}），此外还有动量为 $-{k_a}/{2 b}$ 的标准穿刺点。在自由费米子的语境中，与非正则奇点相联系的算符最初是在文献~\cite{Miwa,Moore} 中定义的。
